% Style based on the supplied Waseem_Generic_AI_ML_GenAI_AgenticAI.tex.
% Python replaces template tokens. Edit layout here, content in JSON/portfolio JS.
\documentclass[9pt,letterpaper]{extarticle}
\usepackage[left=0.68in,right=0.68in,top=0.45in,bottom=0.45in]{geometry}
\usepackage[T1]{fontenc}
\input{glyphtounicode}
\pdfgentounicode=1
\usepackage[utf8]{inputenc}
\usepackage{lmodern}
\usepackage{microtype}
\usepackage{pifont}
\usepackage{xcolor}
\usepackage{enumitem}
\usepackage{titlesec}
\usepackage[hidelinks]{hyperref}
\usepackage{tabularx,array,parskip,ragged2e,needspace}
\pagestyle{empty}
\setlength{\parindent}{0pt}
\setlength{\parskip}{1pt}
\definecolor{accent}{HTML}{155B8A}
\titleformat{\section}{\large\bfseries\color{accent}}{}{0pt}{}[\vspace{0.5pt}\color{accent}\titlerule]
\titlespacing*{\section}{0pt}{3pt}{2pt}
% Compact subsections beneath the main accent-rule sections.
\titleformat{\subsection}{\normalsize\bfseries\color{accent}}{}{0pt}{}
\titlespacing*{\subsection}{0pt}{3pt}{1pt}
\setlist[itemize]{leftmargin=1.15em,itemsep=0.8pt,topsep=1pt,parsep=0pt,partopsep=0pt}
\newcommand{\role}[4]{%
  \Needspace{4\baselineskip}\textbf{#1}\hfill\textbf{#2}\\[-1pt]
  \textit{#3}\hfill\textit{#4}\par}
\newcommand{\entry}[2]{\Needspace{4\baselineskip}\textbf{#1}\hfill\textit{#2}\par}
\newcommand{\skillrow}[2]{\textbf{#1:} & #2\\}
\hypersetup{pdftitle={Waseem Raza, PhD - Generative \& Agentic AI},pdfauthor={Waseem Raza, PhD}}
\begin{document}

\begin{center}
{\LARGE\bfseries Waseem Raza, PhD}{\normalsize\mdseries }\\[2pt]
{\large\bfseries Machine Learning Engineer \textbar{} Generative \& Agentic AI}\\[1pt]
{\small LLM Applications \textbar{} RAG \textbar{} Agent Orchestration \textbar{} Cloud Data Pipelines}\\[3pt]
{\footnotesize \ding{41}\enspace \href{mailto:waseem.placements@gmail.com}{waseem.placements@gmail.com} \textbar{} \href{mailto:waseemraza844@gmail.com}{waseemraza844@gmail.com}\\[1pt]New Jersey, United States \quad \ding{37}\enspace \href{tel:+17328039703}{+1 732-803-9703}\\[1pt]\href{https://www.linkedin.com/in/waseem-raza-phd-9a737858/}{LinkedIn} \textbar{} \href{https://github.com/WaseemRaza844}{GitHub} \textbar{} \href{https://scholar.google.com/citations?user=d58Drr0AAAAJ\&hl=en}{Google Scholar} \textbar{} \href{https://waseemraza844.github.io/portfolio/learning.html}{Coursera Courses} \textbar{} \href{https://waseemraza844.github.io/portfolio/}{Portfolio}}
\end{center}
\vspace{-5pt}
\Needspace{3\baselineskip}
\section{Professional Profile}
Generative AI and Machine Learning Engineer with a PhD in Electrical and Computer Engineering and 7+ years of combined industry and research experience. Build LLM applications, retrieval-augmented generation (RAG) pipelines, structured information-extraction workflows, tool-using agents, and AI-ready data pipelines using Python, LangChain, LangGraph, Hugging Face, OpenAI APIs, Azure Databricks, and cloud services. Experience includes multimodal diagnostic prototypes, prompt engineering, embeddings, semantic retrieval, agent orchestration, model evaluation, domain adaptation, and production-oriented AI delivery.
\Needspace{2\baselineskip}
\subsection{Technical Competencies}
\Needspace{2\baselineskip}\noindent{\small\textbf{Generative \& Agentic AI:} LLM applications, retrieval-augmented generation (RAG), prompt engineering, document ingestion, chunking, embeddings, semantic retrieval, structured information extraction, tool-using agents, agent orchestration, and grounded response generation}\par
\Needspace{2\baselineskip}\noindent{\small\textbf{Microsoft AI Ecosystem:} Microsoft Foundry, Azure OpenAI, Azure AI Services, Azure AI Search, Azure Machine Learning, Foundry models and agents, Azure Content Understanding, Semantic Kernel, and responsible AI concepts}\par
\Needspace{2\baselineskip}\noindent{\small\textbf{LLM Frameworks \& Models:} Hugging Face Transformers, LangChain, LangGraph, OpenAI APIs, transformer architectures, fine-tuning concepts, NLP foundation models, prompt templates, structured outputs, and LLM evaluation}\par
\Needspace{2\baselineskip}\noindent{\small\textbf{Data \& AI Platforms:} Azure Databricks, Apache Spark, PySpark, Azure Data Factory, Azure Synapse Analytics, Azure Functions, Azure Blob Storage, Azure Data Lake, Google BigQuery, and cloud analytics workflows}\par
\Needspace{2\baselineskip}\noindent{\small\textbf{Data Engineering:} ETL/ELT pipelines, AI-ready datasets, batch processing, ingestion, cleansing, schema transformation, validation, feature engineering, data-quality assessment, and structured and unstructured data workflows}\par
\Needspace{2\baselineskip}\noindent{\small\textbf{Machine Learning:} PyTorch, TensorFlow, scikit-learn, XGBoost, transfer learning, domain adaptation, adversarial learning, anomaly detection, classification, regression, time-series forecasting, cross-validation, and model evaluation}\par
\Needspace{2\baselineskip}\noindent{\small\textbf{MLOps \& Software Engineering:} Docker, Kubernetes, OpenShift, Git, GitHub Actions, CI/CD, REST APIs, Flask, testing, Linux, shell scripting, model deployment, reproducible workflows, and cloud-native application development}\par
\Needspace{2\baselineskip}\noindent{\small\textbf{Programming \& Analytics:} Python, SQL, C/C++, and R; Pandas, NumPy, SciPy, Statsmodels, Matplotlib, Plotly, Streamlit, Jupyter, and exploratory data analysis}\par
\Needspace{4\baselineskip}\textbf{Leadership and Professional Strengths}\par
\begin{itemize}
\item \textbf{Ownership and accountability:} Lead AI/ML initiatives through requirements analysis, data preparation, model development, validation, stakeholder review, and deployment planning.
\item \textbf{Technical leadership:} Define technical approaches, evaluate alternatives, coordinate experimentation, and communicate findings to guide projects from ambiguous requirements to measurable outcomes.
\item \textbf{Cross-functional collaboration:} Work with software engineers, network specialists, researchers, geologists, and public-safety stakeholders to translate domain requirements into practical data and AI solutions.
\item \textbf{Continuous learning:} Apply ongoing study in generative AI, agentic AI, RAG, LLM evaluation, and cloud-native platforms through hands-on projects and professional coursework.
\item \textbf{Reliability and commitment:} Maintain reproducible workflows, careful validation, and attention to detail throughout experimentation and delivery.
\item \textbf{Adaptability and problem solving:} Combine critical thinking, systematic experimentation, and data-driven analysis to address unfamiliar technologies and complex operational problems.
\end{itemize}
\Needspace{3\baselineskip}
\section{Professional Experience}
\role{Machine Learning Engineer}{02/2025 - Present}{AT\&T, contractor via Innovatix Inc.}{Bedminster, NJ}
\begin{itemize}\interlinepenalty=10000
\item Lead end-to-end development and validation of AI-driven E911 diagnostic workflows, combining NLP, domain-adaptive machine learning, and production-oriented data pipelines for critical service assurance.
\item Develop multimodal GenAI-assisted diagnostic prototypes that combine work-order text, location information, network context, and operational signals for VoNR and 5G service-assurance use cases.
\item Design reusable model-validation, error-analysis, and failure-mode assessment workflows to compare performance across heterogeneous operational domains and distribution shifts.
\item Partner with network engineers, software teams, and public-safety stakeholders to translate operational requirements into validated AI workflows and deployable data assets.
\end{itemize}
\role{Data Scientist and AI/ML Engineer}{09/2024 - 12/2024}{Impac Exploration Services}{Weatherford, OK}
\begin{itemize}\interlinepenalty=10000
\item Built an LLM-enabled analytical pipeline for extracting and interpreting information from geochemical well data using Python and Azure-based data services.
\item Developed scalable ingestion, cleansing, transformation, and validation workflows using Azure Databricks, Azure Functions, and Synapse Analytics.
\item Automated ETL processing for large exploration datasets, reducing data-preparation time by 40\% and improving reproducibility.
\item Collaborated with geologists to translate domain requirements into AI-ready datasets and structured analytical outputs.
\end{itemize}
\role{Advanced-Degrees Intern}{06/2024 - 08/2024}{AT\&T Labs}{Austin, TX}
\begin{itemize}\interlinepenalty=10000
\item Analyzed 500K+ LTE/5G-NR diagnostic logs and applied ML methods to detect anomalies, reducing diagnostic time by 20\%.
\item Investigated 5G NR-NTN cell-selection behavior and improved decision algorithms by 15\% through data-driven analysis and validation.
\item Conducted structured analysis of large operational datasets and communicated technical findings to engineering teams supporting next-generation wireless systems.
\end{itemize}
\role{Advanced Wireless Technology Development Intern}{09/2022 - 01/2023}{Futurewei Technologies}{Schaumburg, IL}
\begin{itemize}\interlinepenalty=10000
\item Improved autoencoder-based channel-state-information compression by 25\% using transfer learning and systematic experimentation on large wireless datasets.
\item Conducted 50+ experiments to validate model performance, robustness, and generalization under changing data conditions.
\item Collaborated with research and engineering stakeholders to translate modeling results into recommendations for next-generation communication systems.
\end{itemize}
\role{Advanced-Degrees Intern}{06/2022 - 08/2022}{AT\&T Labs}{Bedminster, NJ}
\begin{itemize}\interlinepenalty=10000
\item Designed ML and statistical methods for analyzing users' timing in 5G standalone networks, improving analytical precision by 20\%.
\item Analyzed large-scale user trace data with SQL and visualization tools and communicated findings to senior technical leadership.
\item Converted raw operational traces into interpretable performance indicators to support diagnosis and network optimization.
\end{itemize}
\role{Graduate Research Assistant}{09/2019 - 07/2025}{AI4Networks, University of Oklahoma}{Tulsa, OK}
\begin{itemize}\interlinepenalty=10000
\item Designed reusable Python simulation and data-generation platforms for large-scale experimentation, including a 3GPP-compliant network simulator and an AI-enabled 5G testbed.
\item Built transfer-learning, adversarial-learning, GAN, autoencoder, and reinforcement-learning pipelines for sparse-data and cross-domain modeling.
\item Developed robustness evaluation frameworks for models operating under distribution shifts, missing data, and limited labeled samples; led 100+ experiments from problem formulation through validation.
\item Developed scalable self-healing and outage-management frameworks that improved modeled network resilience by 30\%.
\item Published peer-reviewed research, presented technical results internationally, and mentored junior researchers in ML implementation, experimental design, and technical writing.
\end{itemize}
\Needspace{3\baselineskip}
\section{Education}
\entry{PhD, Electrical and Computer Engineering}{09/2019 - 07/2025}
University of Oklahoma, Norman, Oklahoma, USA\par
\entry{MSc, Telecommunication Engineering}{08/2014 - 08/2016}
University of Engineering and Technology, Taxila, Pakistan\par
\entry{BSc, Telecommunication Engineering}{11/2010 - 06/2014}
University of Engineering and Technology, Taxila, Pakistan\par
\Needspace{3\baselineskip}
\section{Selected Projects \& Research}
\Needspace{2\baselineskip}
\subsection{Projects \& Applied Research}
\entry{\href{https://github.com/WaseemRaza844/agentic-ai-projects-by-was}{Agentic AI and RAG Applications}}{Portfolio Project \textbar{} 2026}
Built grounded GenAI workflows that combine document retrieval, vector search, stateful orchestration, tool-using agents, memory, evaluation, and multi-agent handoffs. Emphasized controllable execution, traceable context flow, and reliable responses for multi-step tasks rather than single unconstrained model calls.\par
{\small\textit{Tools \& Methods:} Python, LangGraph, LangChain, CrewAI, AutoGen, RAG, Vector Databases, Embeddings, Tool Calling, Agent Memory, Evaluation\par}
\vspace{1pt}
\entry{Multi-Agent Orchestration with LangGraph, CrewAI \& AutoGen}{Certification Project \textbar{} 2026}
Implemented stateful multi-agent workflows with tool calling, memory, conditional routing, ReAct-style reasoning, specialized-agent handoffs, and synthesis across LangGraph, CrewAI, AutoGen/AG2, and BeeAI. Focused on explicit orchestration patterns that make multi-agent behavior controllable and reusable.\par
{\small\textit{Tools \& Methods:} LangGraph, CrewAI, AutoGen / AG2, BeeAI, Tool Calling, Agent Memory, Conditional Routing, ReAct, Multi-Agent Systems\par}
\vspace{1pt}
\entry{RAG Retrieval Stack with ChromaDB, FAISS \& LlamaIndex}{Certification Project \textbar{} 2026}
Built retrieval workflows covering document processing, chunking, embeddings, vector indexing, semantic search, advanced retrievers, and interactive RAG application flows. Compared retrieval components and connected retrieved context to LLM generation through LangChain and LlamaIndex orchestration.\par
{\small\textit{Tools \& Methods:} RAG, ChromaDB, FAISS, LlamaIndex, LangChain, Embeddings, HNSW, Semantic Search, Retriever Design, Gradio\par}
\vspace{1pt}
\entry{Amazon Bedrock Generative AI \& Agent Workflows}{Certification Project \textbar{} 2026}
Implemented foundation-model prompting, knowledge-base retrieval, RAG integration, model-customization exercises, workflow automation, and intelligent-agent patterns using Amazon Bedrock. Explored production-oriented cloud GenAI patterns with managed model access, retrieval-backed grounding, and reusable application workflows.\par
{\small\textit{Tools \& Methods:} Amazon Bedrock, Foundation Models, RAG, Knowledge Bases, LangChain, AI Agents, Prompt Engineering, AWS GenAI Services, Workflow Automation\par}
\vspace{1pt}
\entry{OCR \& Document Intelligence Extraction Pipeline}{Guided Project \textbar{} 2025}
Implemented a document-intelligence pipeline that converts reports, invoices, and contract-style documents into structured outputs through OCR, layout/table parsing, entity and field extraction, normalization, and validation. Designed the workflow for downstream analytics and data integration rather than isolated text extraction.\par
{\small\textit{Tools \& Methods:} OCR, Document Intelligence, LLMs, Named Entity Recognition, Table Extraction, Structured JSON, Normalization, Validation, Data Integration\par}
\vspace{1pt}
\entry{Automated Audio Meeting Minutes with Whisper \& GPT}{Guided Project \textbar{} 2025}
Built a multimodal meeting workflow that converts recorded audio to text, cleans and chunks transcripts, generates structured summaries, and extracts decisions, owners, and action items. Combined speech recognition with LLM-based structured generation to automate repeatable knowledge-capture tasks.\par
{\small\textit{Tools \& Methods:} Python, OpenAI Whisper, Speech-to-Text, GPT / LLMs, Prompt Design, Structured Output, Summarization, Workflow Automation, Multimodal AI\par}
\vspace{1pt}
\entry{\href{https://github.com/WaseemRaza844/system-design-masterclass}{System Design Case-Study Portfolio}}{Guided Project \textbar{} 2026}
Developed architecture case studies that translate ambiguous product requirements and traffic assumptions into service boundaries, APIs, data models, storage choices, caching strategies, and reliability targets. Documented scalability, consistency, capacity, and failure-mode tradeoffs to support production-oriented system design decisions.\par
{\small\textit{Tools \& Methods:} System Design, REST APIs, Distributed Systems, Data Modeling, Caching, Scalability, Reliability, Capacity Planning, Architecture Tradeoffs\par}
\vspace{1pt}
\Needspace{2\baselineskip}
\subsection{Publications}
\entry{Real-World Medication Adherence Monitoring Using Hybrid Neural Networks and Smartwatches}{2026}
\textit{HealthINF \textbar{} Accepted}\par
\entry{\href{https://doi.org/10.3390/s25247651}{On Multi-Parameter Optimization and Proactive Reliability in 5G and Beyond Cellular Networks}}{2025}
\textit{Sensors 2025}\par
\entry{\href{https://doi.org/10.1109/PIMRC59610.2024.10817403}{An AI-Driven Framework for Enhancing Resilience in Propagation Models to Enable Digital Twin}}{2024}
\textit{IEEE PIMRC}\par
\Needspace{3\baselineskip}
\section{Certifications \& Professional Development}
\Needspace{3\baselineskip}\noindent\textbf{1. Generative AI Applications with Amazon Bedrock} {\small --- \textcolor{blue}{\underline{\href{https://www.coursera.org/account/accomplishments/verify/3W16P8EE4BJZ}{Coursera}}} \textbar{} Completed \textbar{} August 27, 2026 \textbar{} \textcolor{blue}{\underline{\href{https://waseemraza844.github.io/portfolio/certificates/genai/aws-bedrock/c2-generative-ai-applications-bedrock.pdf}{Certificate}}}}\par
\textbf{Focus:} Covers development of generative-AI applications on Amazon Bedrock, including foundation-model access, prompt-based application logic, model interaction, and managed AWS services. Emphasizes practical cloud patterns for building and operating LLM-powered applications.\par
{\small\textbf{Tools \& Skills:} Amazon Bedrock, foundation models, prompt engineering, AWS generative-AI services, LLM application development, cloud AI workflows\par}
\vspace{2pt}
\Needspace{3\baselineskip}\noindent\textbf{2. Generative AI Engineering and Fine-Tuning Transformers} {\small --- \textcolor{blue}{\underline{\href{https://www.coursera.org/account/accomplishments/verify/SG4GDJ4CLPVL}{IBM / Coursera}}} \textbar{} Completed \textbar{} August 6, 2026}\par
\textbf{Focus:} Covers transformer-based generative-AI engineering, including pretrained language models, model adaptation, fine-tuning strategies, and evaluation. Emphasizes converting reusable foundation models into task-focused systems while understanding training and inference tradeoffs.\par
{\small\textbf{Tools \& Skills:} transformers, foundation models, fine-tuning, transfer learning, tokenization concepts, model adaptation, LLM evaluation\par}
\vspace{2pt}
\Needspace{3\baselineskip}\noindent\textbf{3. Agentic AI with LangChain and LangGraph} {\small --- \textcolor{blue}{\underline{\href{https://www.coursera.org/account/accomplishments/verify/L4EY4SLGCPUE}{Coursera}}} \textbar{} Completed \textbar{} February 9, 2026}\par
\textbf{Focus:} Covers design of stateful agentic applications using LangChain and LangGraph, including graph-based orchestration, tool use, memory and state, routing, and multi-step workflows. Emphasizes controllable agent execution and composition of reliable LLM-driven processes.\par
{\small\textbf{Tools \& Skills:} LangChain, LangGraph, tool-calling, state and memory, workflow orchestration, routing, agent design\par}
\vspace{2pt}
\Needspace{3\baselineskip}\noindent\textbf{4. Generative AI: Foundation Models and Platforms} {\small --- \textcolor{blue}{\underline{\href{https://www.coursera.org/account/accomplishments/verify/AXQJZ5PZ3WH5}{IBM / Coursera}}} \textbar{} Completed \textbar{} October 1, 2024}\par
\textbf{Focus:} Introduces foundation-model concepts and the platforms used to build generative-AI applications, including model capabilities, model selection, prompting, and application integration. Emphasizes how pretrained models serve as reusable components across modern GenAI workflows.\par
{\small\textbf{Tools \& Skills:} foundation models, LLM concepts, model selection, prompt engineering, generative-AI platforms, application integration, responsible model use\par}
\vspace{2pt}
\Needspace{3\baselineskip}\noindent\textbf{5. Advanced RAG with Vector Databases and Retrievers} {\small --- \textcolor{blue}{\underline{\href{https://www.coursera.org/account/accomplishments/verify/AH57KRK80TBQ}{Coursera}}} \textbar{} Completed \textbar{} January 24, 2026}\par
\textbf{Focus:} Covers advanced retrieval-augmented generation pipelines that improve grounding through vector databases, retriever design, indexing, semantic search, and retrieval strategies. Emphasizes selecting and evaluating retrieval components that provide more relevant context to LLM applications.\par
{\small\textbf{Tools \& Skills:} RAG, vector databases, embeddings, semantic retrieval, retriever design, chunking and indexing, retrieval evaluation\par}
\vspace{2pt}
\Needspace{3\baselineskip}\noindent\textbf{6. Build AI Agents using MCP} {\small --- \textcolor{blue}{\underline{\href{https://www.coursera.org/account/accomplishments/verify/QWAQ3NL4TGSZ}{IBM / Coursera}}} \textbar{} Completed \textbar{} June 28, 2026}\par
\textbf{Focus:} Covers construction of AI agents that use Model Context Protocol to connect language models with external tools, services, and structured context. Emphasizes interoperable tool interfaces, context-aware actions, and agent workflows that combine reasoning with external capabilities.\par
{\small\textbf{Tools \& Skills:} Model Context Protocol, tool-using agents, context integration, LLM orchestration, API and tool interfaces, agent workflow design\par}
\vspace{2pt}
\end{document}
